% BAN TIENG VIET cua 06_discussion_conclusion.tex (2026-10-08). So lieu giu nguyen.
\section{Bàn luận}
\label{sec:discussion}

\subsection{Quan hệ với các công trình trước}

Việc giám sát AWD đã được tiếp cận từ ba hướng --- cảm biến mực nước có thiết
bị~\cite{delacruz2022}, viễn thám~\cite{lovell2019,islam2026alose,mlrice2025},
và MRV ở cấp chương trình~\cite{qd1490,mae2026progress,irri2026mrv} --- trong
khi phân tích tổng hợp nông học~\cite{carrijo2017} và các nghiên cứu về tiếp
nhận~\cite{lampayan2015,enriquez2021} giải thích vì sao giám sát có thiết bị
không lan rộng được. Không hướng nào đánh giá \emph{độ tin cậy của chính bản
ghi hành động}, và đó là câu hỏi được đặt ra ở đây (khảo sát đầy đủ ở
Phụ lục~S8). Với radar vệ tinh, phương án chi phí thấp hiển nhiên, chúng tôi đo
chứ không tranh luận: Bảng~\ref{tab:benchmark} chấm điểm mọi phương án theo số
pha khô đủ điều kiện cấp tín chỉ mà nó xác nhận được trong $150$ pha, một pha
chỉ được tính nếu có ít nhất một lần quan sát rơi vào trong
nó~\cite{clauss2017s1,hashemi2022s1,delacruz2022}. Hai điểm cần thận trọng: cột
này đo khả năng \emph{kiếm} tín chỉ chứ không đo khả năng phát hiện gian lận ---
ống đo và bản ghi lệnh cùng đạt $100\%$ nhưng khác nhau về chi phí và về thứ
chúng chứng minh được --- và ảnh chụp kèm tự khai đạt $0\%$ vì nó không tạo ra
chuỗi quan sát nào để đối chiếu. Do đó radar vẫn giữ vai trò mà Phụ lục 4 của
VM0051 giao cho nó: một phép kiểm chéo theo thửa làm giảm tần suất kiểm tra thực
địa~\cite{vm0051} (Phụ lục~S9.2).

\subsection{Giới hạn}

Bảy giới hạn ràng buộc các tuyên bố, tất cả đều được phát biểu đầy đủ ở
Phụ lục~S21; hai giới hạn chi phối các kết quả chính là như sau. \emph{Nông dân
mô phỏng}: các bản giả mạo được tiêm vào quỹ đạo mô phỏng, nên vẫn cần một thử
nghiệm thực địa để đo phản ứng hành vi --- trên hết là liệu nông dân có thay đổi
việc bơm \emph{thật sự} của mình một khi biết có neo thiết bị hay không, một tấn
công làm thay đổi $M$ chứ không phải $L$ và mô hình của chúng tôi không biểu
diễn được. \emph{Tấn công chỉ được minh hoạ trong năm khan nước}: năm 2025 phần
lớn nhật ký rơi xuống dưới mức tối thiểu ba lệnh mà tại đó một tấn công thời điểm
mới có thể dựng được --- trong $37$ nhật ký thật sự bị sửa ở $\Theta$ đầy đủ, cả
$37$ đều thuộc năm 2024. Chính vách ngăn giữa hai năm đó là một phát hiện (rủi
ro gian lận MRV tập trung ở nơi khan nước), nhưng kết quả tấn công chưa phải là
bằng chứng về một mùa mưa; ở $\Theta$ rút gọn có mười hai nhật ký bị sửa thuộc
năm 2025 và chúng hành xử giống nhau --- gợi ý chứ chưa kết luận. Năm giới hạn
còn lại được liệt kê ở S21.

\section{Kết luận}

Chúng tôi khởi đầu với mục tiêu làm cho MRV lúa gần như miễn phí bằng cách kiểm
toán một nhật ký tưới do nông dân tự ghi đối chiếu với vật lý nước trong đất.
Phiên bản đầu tiên của lớp đó báo cáo không có kết tội oan và phát hiện trọn
vẹn; cả hai con số đều là sản phẩm của một thửa ruộng do chính mô hình kiểm
toán sinh ra, và việc dựng lại thửa ruộng đó trước một đối thủ tối ưu hoá đại
lượng được cấp tín chỉ đã làm thay đổi kết luận.

Bộ phát hiện mô hình điểm kết tội oan $26{,}7\%$ nông dân trung thực; một ngưỡng
tỉ lệ và một phép giao trên tập bất định $150$ phần tử hạ con số đó xuống
$5{,}0\%$ và $0/60$, với cái giá đo được về năng lực phát hiện ($9/9$ xuống
$4/9$ ở gian lận bỏ sót). Trong khi đó một phép hoán vị thời điểm khai báo giữ
nguyên tổng lượng nước làm thổi phồng số pha khô được cấp tín chỉ từ $33\%$ đến
$38\%$, đồng thời lọt qua tầng bền vững ở $8$ trong $9$ và $3$ trong $3$ nhật ký
bị sửa, và không cấu hình vật lý nào thoát khỏi sự đánh đổi đó: hoãn một khai báo
vừa kéo dài pha khô vừa làm tăng phần thiếu hụt dành cho nó, nên lợi nhuận và sự
không thể phát hiện dùng chung một gradient.

Cách sửa không thuộc về thuỷ văn. Đối chiếu nhật ký khai báo với bản ghi lệnh có
chữ ký và dấu thời gian của gateway bắt được $116/116$ nhật ký thật sự bị giả
mạo và không kết tội oan nhật ký nào trong $60$ nhật ký trung thực, với chi phí
không vượt quá một bộ điều khiển mà công trình đồng hành đã triển khai. Điều còn
lại của tiền đề ban đầu là: vật lý \emph{ràng buộc} một nhật ký, nhưng chỉ một
bản ghi thiết bị có chữ ký mới \emph{chứng thực} được nó --- MRV chi phí gần như
bằng không chỉ khả thi ở nơi có một thiết bị đứng giữa nông dân và máy bơm. Ba
con số đã công bố trước đó được sửa lại tương ứng ($[49{,}0; 86{,}0]$~mm thay vì
$61$~mm; một bộ lọc thứ ba bất động; một phép thử chênh lệch tuyệt đối không thể
chấp nhận). Mã nguồn, cả hai tập bất định, toàn bộ $281$ nhật ký đã kiểm toán và
$46$ phép thử đơn vị đều được công khai.
